\documentclass[UTF8,11pt,a4paper]{ctexart}

\usepackage[margin=2.1cm,headheight=15pt]{geometry}
\usepackage{amsmath,amssymb}
\usepackage{booktabs}
\usepackage{longtable}
\usepackage{array}
\usepackage{tabularx}
\usepackage{enumitem}
\usepackage{xcolor}
\usepackage{tcolorbox}
\usepackage{fancyhdr}
\usepackage{hyperref}
\hypersetup{colorlinks=true,linkcolor=blue!45!black,urlcolor=blue!45!black}

\definecolor{ink}{HTML}{1F2937}
\definecolor{accent}{HTML}{B45309}
\definecolor{muted}{HTML}{6B7280}

\ctexset{
  section    = {format = {\Large\bfseries\color{accent}}},
  subsection = {format = {\large\bfseries\color{ink}}},
}

\pagestyle{fancy}
\fancyhf{}
\fancyhead[L]{\small\color{muted}SOE104 CW2 · Oregon Trail}
\fancyhead[R]{\small\color{muted}决策模型 · v1.1 · 组内阅读版}
\fancyfoot[C]{\small\color{muted}\thepage}
\renewcommand{\headrulewidth}{0.4pt}

\setlist{itemsep=2pt,topsep=3pt,parsep=0pt}
\setlength{\parskip}{0.35em}
\renewcommand{\arraystretch}{1.25}

\newtcolorbox{keybox}[1][]{
  colback=orange!4, colframe=orange!55!black, boxrule=0.6pt, arc=2pt,
  left=7pt, right=7pt, top=5pt, bottom=5pt, fonttitle=\bfseries, #1
}
\newtcolorbox{warnbox}[1][]{
  colback=red!3, colframe=red!45!black, boxrule=0.6pt, arc=2pt,
  left=7pt, right=7pt, top=5pt, bottom=5pt, fonttitle=\bfseries, #1
}
\newtcolorbox{goodbox}[1][]{
  colback=green!3, colframe=green!35!black, boxrule=0.6pt, arc=2pt,
  left=7pt, right=7pt, top=5pt, bottom=5pt, fonttitle=\bfseries, #1
}

\title{\vspace{-1.2cm}\textbf{CW2 决策模型}\\[4pt]
\large 如何把“最优解”推导出来，而不是猜出来\\[2pt]
\normalsize 依据 Engine v1.1 / Marking v1.1}
\author{Group A--E\\\small 整理：A}
\date{2026-11-22 \quad v2.1}

\begin{document}
\maketitle
\thispagestyle{fancy}

\begin{keybox}[title={这份文档要解决什么}]
前三份讲“规则”“评分”“要做什么”。这一份只讲一件事：
\textbf{给定状态和严重度，怎么算出一个能被论证的选择。}

核心主张不变：\textbf{不要背结论表，要背公式。}但 v1.1 改了核心数值，
所以\underline{所有旧公式的系数都要重算}——本文就是重算的结果。
\end{keybox}

\begin{warnbox}[title={相对 v1.0 文档的三处更正}]
\textbf{1. v1.0 里“SAFE 代价不随 $s$ 变化，所以高 $s$ 时 SAFE 反超”的结论作废。}
v1.1 的 SAFE 已经随 $s$ 缩放，不能再当成常数。

\textbf{2. 影子价格全部变了。}食物变贵（$\lambda_f$ 从 1.25 降到 1.0）、
HP 变贵（$\lambda_h$ 从 0.8 升到 2.0）、钱变得只在旅途中有效。

\textbf{3. 但有一个结论完全不变：}ILLNESS 的 RISKY/AVOID 比值恒为 $(s+5)/10$，
分界仍然在 $s=5$。这是本文最稳的一条。
\end{warnbox}

%=====================================================================
\section{问题的形式化}

\subsection{目标函数（v1.1）【U】}

\[
\text{Game} = 40\times\text{arrival} + 30\times\text{condition} + 30\times\text{cargo}
\]

\[
\text{arrival}=\frac{D}{2000},\qquad
\text{condition}=\frac{\sum_{i\in\text{初始成员}} h_i}{100\,N},\qquad
\text{cargo}=\min\!\Big(1,\frac{F+O+A+P}{1000}\Big)
\]

其中 $h_i$ 是到达时该成员血量（死者为 0），$N$ 是初始人数。
cargo 里的 $F,O,A,P$ 分别是食物、牛、弹药、零件，\textbf{不含钱}。

\begin{keybox}[title={三个结构性变化}]
\begin{itemize}[leftmargin=1.4em]
  \item \textbf{死亡从“少一个分母”变成“扣 $1/N$ 的状态分”。}
        4 人队死 1 个，condition 上限立刻掉到 0.75。
  \item \textbf{钱彻底退出计分。}它只在旅途中作为购买力存在，终局价值为 0。
  \item \textbf{受伤也扣分。}旧版只要活着就是满分，新版残血到达一样掉 condition。
\end{itemize}
\end{keybox}

\subsection{展开与化简：把 FINAL 变成一把尺子}

\noindent 上面的公式还不能直接用来做决定——它写的是\textbf{终局状态}（走了多远、剩多少血、剩多少货），
而我们要决定的是\textbf{每回合做什么}。中间差一步替换。这一节把那一步补上。

\subsubsection*{第一步：拆出静态项}

\[
\text{FINAL}=0.20\,\text{OOP}+0.50\,\text{Game}+0.15\,\text{Naming}+0.15\,\text{Efficiency}
\]

前三项里，OOP、Naming、Efficiency 由代码写法决定，和跑得好不好无关。记
\[
S=0.20\,\text{OOP}+0.15\,\text{Naming}+0.15\,\text{Efficiency}\in[0,50],
\]
它是个\textbf{常数}——写好代码就锁定了。剩下的全部问题都是让 $0.5\,\text{Game}$ 尽量大：

\[
\text{FINAL}=S+\tfrac12\,\text{Game}.
\]

\begin{keybox}[title={第一个结论}]
\textbf{策略能影响的只有一半的分数。}
把 Game 从 0 拉到满分，最多值 50 分；而静态分写干净就是 50 分，
且不需要任何回放就能拿到。这就是为什么文档反复说“先拿确定的分”。
\end{keybox}

\subsubsection*{第二步：把 Game 拆成三个可加项}

\[
\tfrac12\text{Game}=20\cdot\frac{D}{2000}+15\cdot\frac{\sum_i h_i}{100N}+15\cdot\min\!\Big(1,\frac{U}{1000}\Big)
\]

把系数并进去：

\[
\boxed{\ \text{FINAL}=S+\frac{D}{100}+\frac{0.15}{N}\sum_i h_i+15\min\!\Big(1,\frac{U}{1000}\Big)\ }
\]

其中 $D$ 是终局里程，$h_i$ 是第 $i$ 个初始成员的血量（死者记 0），$U=F+O+A+P$ 是终局货物总量，$N$ 是初始人数。

\subsubsection*{第三步：求边际，得到三种资源的汇率}

对每个变量求偏导，就是“多一单位值多少分”：

\begin{center}
\begin{tabular}{@{}lccl@{}}
\toprule
\textbf{资源} & \textbf{边际分} & \textbf{折成英里} & \textbf{说明} \\
\midrule
1 英里 & $1/100=0.01$ & $1$ & 基准单位 \\
1 点血量（某一人） & $0.15/N$ & $15/N$ & 随人数变化：$N{=}3{:}5.0$，$N{=}4{:}3.75$，$N{=}5{:}3.0$ \\
1 单位货物 & $0.015$ & $1.5$ & 未达 1000 上限时 \\
\bottomrule
\end{tabular}
\end{center}

\noindent 换句话说，\textbf{整个游戏只有一种货币：英里}。

\begin{keybox}[title={三个可以直接用于决策的换算}]
\begin{enumerate}[leftmargin=1.6em]
  \item \textbf{1 点血量 $=$ 3.75 英里}（4 人队）。这就是为什么“掉 4 血换 8 英里”的 GRUELING 是亏的。
  \item \textbf{1 单位货物 $=$ 1.5 英里}。剩 200 食物在车上，等于白送 300 英里。
  \item \textbf{死一个成员}（4 人队）$=100$ 点血量 $=375$ 英里。
        这与 \texttt{02} 里“一条命 $\approx375$ 英里”完全对得上——两条路径算出同一个数。
\end{enumerate}
\end{keybox}

\begin{warnbox}[title={人数越少，每个人越值钱}]
血量项的系数是 $15/N$，所以：
\begin{center}
\begin{tabular}{@{}lccc@{}}
\toprule
& $N=3$ & $N=4$ & $N=5$ \\
\midrule
1 点血量 & 5.0 英里 & 3.75 英里 & 3.0 英里 \\
死一个人 & 500 英里 & 375 英里 & 300 英里 \\
\bottomrule
\end{tabular}
\end{center}
\textbf{这是一个反直觉但重要的结论}：小队虽然干活的人少，但\textbf{每个人的命更贵}。
3 人队死一个，等于丢掉 500 英里的距离分——相当于整局四分之一的进度。
\end{warnbox}

\subsubsection*{第四步：把决策变量代进去}

终局状态由每回合的决策累积而成：

\[
D=\sum_t \Delta D_t,\qquad
\sum_i h_i=\sum_i 100-\sum_t \Delta H_t,\qquad
U=U_0-\sum_t C_t+\sum_t P_t
\]

其中 $\Delta D_t$ 是第 $t$ 回合的里程变化（可为负），$\Delta H_t$ 是全场血量净变化，
$C_t$ 是该回合消耗的资源，$P_t$ 是产出的资源（狩猎、交易所得）。

代回上面的式子，就得到\textbf{每回合的评分式}：

\[
\boxed{\ \Delta\text{FINAL}_t=0.01\,\Delta D_t+\frac{0.15}{N}\,\Delta H_t-0.015\,C_t+0.015\,P_t\ }
\]

\begin{keybox}[title={这一行就是整份文档的总结}]
\textbf{每一个动作、每一个事件姿态，都可以用这一个式子打分：}
它带来多少英里（第一项），带来或损失多少血量（第二项），
消耗多少资源（第三项），产出多少资源（第四项）。

下面几节所有表格里的“英里当量”，都是这个式子的特例。
\end{keybox}

\subsubsection*{第五步：两种价格——为什么 §4 的 $\lambda$ 不是这张表}

上表的数字是\textbf{终局价值}（东西撑到终点值多少）。但在\textbf{旅途中}做决定时，
用的应该是\textbf{替换成本}（现在弄丢一单位，要花多大代价补回来）。两者不同：

\begin{center}
\small
\begin{tabular}{@{}p{2.0cm}p{3.4cm}p{3.4cm}p{5.0cm}@{}}
\toprule
\textbf{资源} & \textbf{终局价值} & \textbf{替换成本（§4）} & \textbf{差在哪} \\
\midrule
食物 & 1.5 英里 & 1.0 英里 & 食物几乎一定会被吃掉，所以按“燃料价”算 \\
血量 & 3.75 英里 & 2.0 英里 & 血可以靠 Doctor 补回来，所以按“补血价”算 \\
弹药 & 1.5 英里 & 1.5 英里 & 只能花钱买，两者相同 \\
\bottomrule
\end{tabular}
\end{center}

\begin{keybox}[title={为什么这个区分很重要}]
\textbf{终局价值 $>$ 替换成本}，说明“补回来”是划算的——差值就是这笔生意的利润。

以血为例：$3.75-2.0=1.75$ 英里/点。只要一次 Doctor HEAL 的 5 点血能活到终点，
按终局口径算值 $5\times3.75=18.75$ 英里，而补回来只花了 2.0 的价格。

\smallskip
\noindent\textbf{但 §8 那张表用的是另一套算法}，不是这张表。同一个 Doctor HEAL：
\[
\underbrace{5\times2.0}_{\text{按替换成本估值}}-\underbrace{3}_{\text{3 食物}}=7.0\ \text{英里}
\]
\textbf{两套口径都对，回答的问题不同}：终局口径问“这 5 点血最终值多少”，
替换成本口径问“这回合赚了多少”。后者更适合逐回合做决定，所以 §8 用它。

反过来说，如果伤害\textbf{补不回来}（回合不够、Doctor 死了、连续挨饿），
就必须按终局价值 3.75 算——这正是 §7.4 里 R1 规则存在的原因。
\end{keybox}

\begin{warnbox}[title={口径不同，绝对值不同，排序也会变}]
同一组动作，两套口径算出来的数值差很多：

\begin{center}
\small
\begin{tabular}{@{}lcc@{}}
\toprule
\textbf{动作} & \textbf{终局口径} & \textbf{替换成本口径（§8）} \\
\midrule
DEFEND（12 点血） & 45.0 & 12.0 \\
Doctor HEAL（5 点血） & 14.25 & 7.0 \\
HUNT（Hunter） & 12.0 & 7.0 \\
REST（2 点血） & 7.5 & 4.0 \\
SCOUT & 6.0 & 6.0 \\
\bottomrule
\end{tabular}
\end{center}

\noindent 差别全部来自血量的单价：终局 3.75/点，替换 2.0/点。所以\textbf{排序会变}：

\noindent 两处名次变化：
\begin{itemize}[leftmargin=1.4em]
  \item \textbf{REST 与 SCOUT 互换}：终局口径下 REST（7.5）压过 SCOUT（6.0）；
        替换口径下反过来（4.0 对 6.0）。
  \item \textbf{HEAL 与 HUNT 在替换口径下并列}（都是 7.0），终局口径下 HEAL 略高（14.25 对 12）。
\end{itemize}
\textbf{规律}：换用终局口径，\textbf{一切“加血”的动作名次上升}，“加里程”的相对下降。

\textbf{实用结论}：用替换成本口径做日常决策；
当\textbf{补血能力饱和}（Doctor 每回合只能救一个人、或连续挨饿）时切到终局口径，
这时“回一点血”比“多走一英里”值钱得多。
\end{warnbox}

\subsection{为什么不能直接求最优}

三个障碍没有变：引擎不在本地（jar 只是桩）、\texttt{Resolver} 看不到车队、随机性。
所以“求最优”仍然是指：\textbf{每一层用可解释的近似，再用回放校准。}

%=====================================================================
\section{v1.1 对模型的六个影响}

\begin{center}
\small
\begin{longtable}{@{}p{4.4cm}p{4.6cm}p{5.6cm}@{}}
\toprule
\textbf{改动} & \textbf{对模型的影响} & \textbf{处理方式} \\
\midrule
\endhead
事件概率 20\%$\to$56\% & 事件次数从约 29 次升到约 38 次 & 事件策略的权重上升 \\
SAFE 随 $s$ 缩放 & SAFE 不再是常数 & 每个事件的 SAFE 都要重新代入 \\
HUNT 2 弹药 & 食物产能腰斩 & $\lambda_a$、$\lambda_f$ 全部重算 \\
HEAL 要 3 食物 & HP 变贵约 2.5 倍 & $\lambda_h$ 从 0.8 升到 2.0 \\
Doctor 光环减伤 & 全队事件伤害 $\times 0.5$ & 引入系数 $d$ \\
口粮 1/3/6 & 食物消耗 +25\% & 食物预算必须重算 \\
\bottomrule
\end{longtable}
\end{center}

\begin{keybox}[title={引入 Doctor 系数 $d$}]
\[
d=\begin{cases}0.5 & \text{有 Doctor 活着}\\ 1.0 & \text{没有}\end{cases}
\]
它乘在\textbf{所有事件的 HP 伤害}上。

重要性质：$d$ 在大多数比较里会\textbf{整体约掉}（见 \S\ref{sec:illness}），
所以它不改变姿态排序，只改变绝对值。唯一例外是与\textbf{固定代价}（死牛）比较时。
\end{keybox}

%=====================================================================
\section{严重度 $s$：能读到，但新公式未公布}

\subsection{已知}

v1.0 的公式是【S】：
\[
s=\operatorname{clamp}\Big(\text{uniform}\{1..5\}+\operatorname{round}(6p),\ 1,\ 10\Big)
\]

Trail Update 1.1 只说“\textbf{Severity now climbs even higher}”，\textbf{没有给新公式}【U】。

\subsection{能确定的两件事}

\begin{enumerate}[leftmargin=1.8em]
  \item \textbf{$s$ 在决策时是已知的。}\texttt{Event.severity()} 在事件发生时返回真实值，
        \textbf{不需要预测}。所有规则都可以直接按 $s$ 分档。
  \item \textbf{后期的 $s$ 更高。}无论新公式是什么，$s$ 随进度上升这一点不变。
\end{enumerate}

\begin{warnbox}[title={待确认}]
新公式影响的是\textbf{分段阈值的位置}（比如 “$s\ge4$ 选 AVOID” 这类），
不影响“按 $s$ 分档”这个方法本身。所以模型可以先上线，阈值后用回放校准。
\end{warnbox}

%=====================================================================
\section{影子价格（v1.1 重算）}

\subsection{定义与基准}

$\lambda_x$ = 一单位资源 $x$ 的边际价值，单位统一用\textbf{英里}。
基准状态：$m=4$ 人、$o=4$ 牛、MEAGER 口粮。

\begin{center}
\small
\begin{longtable}{@{}p{1.6cm}p{1.2cm}p{1.8cm}p{9.8cm}@{}}
\toprule
\textbf{资源} & \textbf{符号} & \textbf{基准值} & \textbf{推导（v1.1）} \\
\midrule
\endhead
食物 & $\lambda_f$ & $\mathbf{1.0}$ & 1 回合 NORMAL $=20$ 英里，MEAGER 吃 $3m+2o=20$ 食物 \\
HP & $\lambda_h$ & $\mathbf{2.0}$ & 最便宜的回血是 Doctor HEAL：动作机会成本 $7.0$ $+$ 食物 $3.0$，换 $5$ HP $\Rightarrow 10/5=2.0$ \\
钱 & $\lambda_m$ & $\mathbf{0.5}$ & 钱只能买食物：$\lambda_m=\lambda_f/P_f$，fair 市场 $P_f=2$ \\
牛 & $\lambda_o$ & $\mathbf{60}$ & \$80 购入，每回合吃 2 食物；跌破 4 头速度 $-25\%$。\textbf{可在交易站补买}，见下方说明 \\
零件 & $\lambda_p$ & $\mathbf{15}$ & \$30 $\times$ 0.5 \\
弹药 & $\lambda_a$ & $\mathbf{1.5}$ & \$3 置换成本；使用价值约 3.5（2 发换 7 食物）\\
\bottomrule
\end{longtable}
\end{center}

\begin{warnbox}[title={$\lambda_h$ 为什么是 2.0（v1.0 文档这里是 0.8）}]
v1.0：HEAL 免费 $+10$，HP 几乎是白送的。
v1.1：回血要花食物和动作。按最便宜的方式（Doctor HEAL）算：
\[
\lambda_h=\frac{\underbrace{\text{动作机会成本}}_{\text{放弃的最好替代}}+\underbrace{3.0}_{\text{3 食物}}}{\underbrace{5}_{\text{HP}}}
\]

\noindent\textbf{“动作机会成本”是什么意思}：每个成员每回合只能做一件事。
他选择 HEAL，就放弃了本可以做的那件事的产出。所以回血的代价不只是 3 食物，
还要加上\textbf{被放弃的那个动作值多少}。

\smallskip
\noindent 取 $7.0$（Hunter 打猎的净值）得到 $\lambda_h=2.0$。\textbf{但这只是区间中点}：

\begin{center}
\small
\begin{tabular}{@{}p{6.4cm}cc@{}}
\toprule
\textbf{不治疗的话这个成员会去做} & \textbf{机会成本} & $\lambda_h$ \\
\midrule
HUNT（Doctor 不是 Hunter，$7-3$） & 4.0 & 1.4 \\
SCOUT（$+6$ 英里） & 6.0 & 1.8 \\
HUNT（按 Hunter 的产出算） & 7.0 & 2.0 \\
DEFEND（受威胁回合，$12\times2.0\times0.5$） & 12.0 & 3.0 \\
\bottomrule
\end{tabular}
\end{center}

\noindent 所以 \textbf{$\lambda_h\in[1.4,\ 3.0]$，取 2.0 是中间值}。
真实值取决于该回合谁去治疗、以及有没有事件——这就是为什么要用回放校准，
而不是把 2.0 当成事实。

\smallskip
\noindent\textbf{校验}：MOVE 用 1 HP 换 5 英里。只要 $\lambda_h<5$，MOVE 就划算——
区间 $[1.4,3.0]$ 全部满足，与“MOVE 值得做”这个观察一致。
若 $\lambda_h$ 超过 5，MOVE 就变成亏本操作，与实际不符，说明估值没跑偏。
\end{warnbox}

\begin{warnbox}[title={$\lambda_o$ 不是永久损失——这里 v1.0 的说法要修正}]
牛可以\textbf{在交易站补买}（TRADER 事件，SAFE 姿态，基础价 \$80，每次最多 50 单位）。
所以「掉一头牛 = 永久出局」是\textbf{过度概括}。

正确的判据是\textbf{回合预算}：
\[
\text{若 }\ \frac{2000-D}{15}>T_{\text{left}}\ \text{则掉第四头牛是致命的}
\]
早期掉牛（比如第 20 回合、还剩 1800 英里 / 100 回合：$1800/15=120>100$）\textbf{必然失败}；
后期掉牛可能还能补回来。所以下面 R2 的规则会写成\textbf{条件式}，而不是一刀切。
\end{warnbox}

\subsection{影子价格随状态变化}

\begin{center}
\small
\begin{longtable}{@{}p{1.8cm}p{5.6cm}p{2.8cm}p{4.4cm}@{}}
\toprule
\textbf{价格} & \textbf{触发条件} & \textbf{调整} & \textbf{直觉} \\
\midrule
\endhead
$\lambda_f$ & 用 BARE 而非 MEAGER & $\times 1.67$ & 同样 20 英里只吃 12 食物 \\
$\lambda_f$ & runway $<3$ & $\times 3\sim5$ & 食物换的是命 \\
$\lambda_f$ & food $>450$ & $\times 0\sim0.3$ & 再多会腐烂 \\
$\lambda_h$ & 任一成员 HP $<25$ & $\times 3\sim\infty$ & 死亡扣 condition \\
$\lambda_h$ & 全队 HP $>75$ 且有 Doctor & $\times 0.6$ & 减伤光环让掉血变慢 \\
$\lambda_m$ & 市场 band $\le0.7$ & $0.71$ & 钱能买便宜食物 \\
$\lambda_m$ & 市场 band $\ge1.4$ & $0.36$ & 钱不值钱 \\
$\lambda_m$ & 进度 $>90\%$ & $\to 0$ & 到终点前花不掉 \\
$\lambda_o$ & oxen $=4$ & $\times 2$ & 再死一头就永远跑不完 \\
$\lambda_o$ & oxen $\ge5$ & $\times 0.6$ & 有缓冲 \\
\bottomrule
\end{longtable}
\end{center}

\begin{keybox}[title={核心算式没变}]
\[
\text{代价}(a)=\sum_x \Delta x_a\cdot\lambda_x,\qquad a^{*}=\arg\min_a \text{代价}(a)
\]
换一种状态，就是换一组 $\lambda$。这就是为什么结论表会失效、公式不会。
\end{keybox}

%=====================================================================
\section{决策面 A：开局采购（v1.1 重算）}

\subsection{食物预算先算清楚}

$m$ 人、$o$ 牛、每回合消耗 $=3m+2o$（MEAGER）。

\begin{center}
\begin{tabular}{@{}lccc@{}}
\toprule
\textbf{配置} & \textbf{每回合} & \textbf{100 回合总需求} & \textbf{起始 500 后缺口} \\
\midrule
4 人 + 5 牛，MEAGER & 22 & 2200 & 1700 \\
4 人 + 4 牛，MEAGER & 20 & 2000 & 1500 \\
4 人 + 4 牛，BARE & 12 & 1200 & 700 \\
3 人 + 4 牛，MEAGER & 17 & 1700 & 1200 \\
3 人 + 4 牛，BARE & 11 & 1100 & 600 \\
5 人 + 4 牛，MEAGER & 23 & 2300 & 1800 \\
\bottomrule
\end{tabular}
\end{center}

\textbf{不动用市场时的供给上限：}

\begin{itemize}[leftmargin=1.5em]
  \item 80 弹药 $\div$ 2 $\times$ 10（Hunter）= \textbf{400 食物}
  \item 剩余 \$380 在便宜市场（\$1.4）= \textbf{271 食物}
  \item 合计约 \textbf{1171}
\end{itemize}

\begin{warnbox}[title={这是 v1.1 最重要的一条结论}]
\textbf{MEAGER 的食物预算在 4 人 4 牛下不闭合}：需要 2000，最多凑到约 1171，缺口 829。

只有两条路能闭合：
\begin{enumerate}[leftmargin=1.5em]
  \item \textbf{降档到 BARE}（需求降到 1200，接近供给上限）；
  \item \textbf{靠 HUNT\_OPP 和市场}（HUNT\_OPP 是 10 食物/弹药，是普通狩猎的 2--3 倍）。
\end{enumerate}
\textbf{所以 v1.1 的默认口粮很可能不是 MEAGER，而是 BARE + 高频治疗。}
这一点必须用回放验证，不能想当然。
\end{warnbox}

\subsection{采购建议}

\begin{center}
\small
\begin{tabular}{@{}lcccccc@{}}
\toprule
\textbf{方案} & \textbf{牛} & \textbf{食物} & \textbf{弹药} & \textbf{零件} & \textbf{剩钱} & \textbf{评价} \\
\midrule
经典 & 5 & 500 & 80 & 2 & \$300 & 第 5 头牛是保险，但每回合多吃 2 \\
省牛 & 4 & 500 & 80 & 2 & \$380 & 4 头刚好满速，风险是掉一头就出局 \\
重弹 & 4 & 500 & 120 & 2 & \$260 & 弹药是食物预算，多买弹等于多买粮 \\
\bottomrule
\end{tabular}
\end{center}

\begin{goodbox}[title={建议}]
\textbf{4 头牛 + 满食物 + 略多弹药}。理由：
\begin{itemize}[leftmargin=1.4em]
  \item 4 头牛是满速的临界点，省下的 \$80 换成 40 单位食物；
  \item 弹药是 v1.1 真正的食物预算（2 发换 10 食物），比直接买食物划算；
  \item 钱不计入 cargo，留着不花是纯损失。
\end{itemize}
\textbf{风险：4 头牛没有缓冲。}一旦在 RIVER 或 BAD\_TRAIL 上死一头，速度掉到 75\%，
2000 英里需要 134 回合——直接出局。所以\textbf{低 $s$ 时可以冒险，高 $s$ 时必须保牛}。
\end{goodbox}

%=====================================================================
\section{决策面 B：配速与口粮（v1.1）}

\subsection{配速}

时间预算不变：NORMAL 在 4 牛下 20 英里/回合，2000 英里需 100 回合，余 20。

\[
\text{needPerTurn}=\frac{2000-D}{120-T}
\]
$\text{needPerTurn}\le20$ 用 NORMAL；否则先让成员 MOVE/SCOUT，仍不够再 GRUELING。

\subsection{GRUELING 的代价变了}

\[
\text{收益}=8\ \text{英里},\qquad
\text{代价}=4m\lambda_h d+0.15\lambda_o=4\times4\times2.0\times0.5+0.15\times60=25\ \text{英里}
\]

\textbf{净亏约 17 英里/回合}（v1.0 是净亏 13.8）。因为 HP 变贵。
\textbf{GRUELING 在 v1.1 里几乎只在最后几回合救急时使用。}

\subsection{口粮的三态（v1.1）}

\begin{center}
\small
\begin{tabular}{@{}lccp{5.0cm}p{4.2cm}@{}}
\toprule
\textbf{口粮} & \textbf{食物/人} & \textbf{HP 变化} & \textbf{折算代价（$m=4$）} & \textbf{何时用} \\
\midrule
BARE & 1 & $-1$/人 & 省 $4\times2\times1=8$ 英里，亏 $4\times1\times2.0=8$ & \textbf{净 0}：只有在食物更紧时才赚 \\
MEAGER & 3 & 0 & 基准 & 默认 \\
FULL & 6 & $+1$/人 & 多花 $4\times3\times1=12$ 英里，赚 $4\times1\times2.0=8$ & 净亏 4：只在濒死时用 \\
\bottomrule
\end{tabular}
\end{center}

\begin{keybox}[title={与 v1.0 相反的结论}]
v1.0 里 FULL 是“$+2$ HP”，净亏 8.6；v1.1 里 FULL 只给 $+1$ HP，净亏 4。

但最关键的变化是 \textbf{BARE 从“净赚 1.8”变成了“净 0”}——
因为 HP 涨价到 2.0，BARE 掉的血正好抵消省下的食物。

\textbf{所以口粮不能靠“划算”来选，要看食物预算：}
\begin{itemize}[leftmargin=1.4em]
  \item 食物够（runway 高）→ \textbf{MEAGER}，血量是免费维持的；
  \item 食物不够（前面算过，这是常态）→ \textbf{BARE}，不是因为它划算，而是因为\textbf{不降档就会饿死}；
  \item 濒死 → \textbf{FULL}，明知亏也要买命。
\end{itemize}
\end{keybox}

%=====================================================================
\section{决策面 C：事件姿态（核心）}\label{sec:events}

\subsection{通用判定式}

\[
C(\text{SAFE})=\sum \Delta_{\text{safe}}\lambda,\qquad
C(\text{RISKY})=\sum_k \Pr(k)\cdot\sum\Delta_k\lambda,\qquad
C(\text{AVOID})=\sum\Delta_{\text{avoid}}\lambda
\]

所有 HP 伤害乘以 Doctor 系数 $d$。

\subsection{基准参数}

$m=4$、$\lambda_f=1.0$、$\lambda_h=2.0$、$\lambda_m=0.5$、$\lambda_o=60$、$\lambda_a=1.5$、$\lambda_p=15$、$d=0.5$。

\noindent 下面反复用到 $m\lambda_h d=4\times2.0\times0.5=4.0$。

\subsubsection*{RIVER}

\begin{align*}
C(\text{SAFE})&=(20+s)\lambda_m=10+0.5s\\
C(\text{RISKY})&=\tfrac{s}{10}\big(10m\lambda_h d+20\lambda_f\big)+0.2\lambda_o=6s+12\\
C(\text{AVOID})&=5
\end{align*}

$s=1$ 时 RISKY 已是 18，$s=10$ 时 72。\textbf{AVOID 恒胜。}

\subsubsection*{STORM —— SAFE 与 RISKY 伤害相同}

\begin{align*}
C(\text{SAFE})&=m s\lambda_h d=4s\\
C(\text{RISKY})&=m s\lambda_h d+0.15\lambda_o=4s+9\\
C(\text{AVOID})&=10
\end{align*}

\textbf{SAFE 严格支配 RISKY}：伤害完全一样，但 RISKY 多 15\% 死牛。
所以只需要比较 SAFE 和 AVOID：$4s<10\Rightarrow s\le2$ 选 SAFE，$s\ge3$ 选 AVOID。

\subsubsection*{BREAKDOWN}

\begin{align*}
C(\text{SAFE})&=\lambda_p=15\quad(\text{有 Mechanic 时}=0)\\
C(\text{RISKY})&=\tfrac{s}{10}\big(5m\lambda_h d+\lambda_p\big)=3.5s\\
C(\text{AVOID})&=15
\end{align*}

有 Mechanic：\textbf{SAFE 免费，恒胜}。
没有 Mechanic：$s\le4$ RISKY（14.0），$s\ge5$ AVOID。

\subsubsection*{ILLNESS —— 唯一不变的结论}\label{sec:illness}

\begin{align*}
C(\text{SAFE})&=\text{官方未公布}\\
C(\text{RISKY})&=\tfrac{s}{10}\cdot m(s+5)\lambda_h d=0.4s(s+5)\\
C(\text{AVOID})&=m s\lambda_h d=4s
\end{align*}

\[
\boxed{\frac{C(\text{RISKY})}{C(\text{AVOID})}=\frac{s+5}{10}}
\]

\textbf{这个比值与 $m$、$\lambda_h$、$d$ 全都无关}——v1.1 改了人数、改了 HP 价格、
加了 Doctor 光环，这个比值\textbf{一点没变}。所以：

\begin{center}
\begin{tabular}{@{}lcccccccccc@{}}
\toprule
$s$ & 1 & 2 & 3 & 4 & \textbf{5} & 6 & 7 & 8 & 9 & 10 \\
\midrule
RISKY & 2.4 & 5.6 & 9.6 & 14.4 & \textbf{20.0} & 26.4 & 33.6 & 41.6 & 50.4 & 60.0 \\
AVOID & 4.0 & 8.0 & 12.0 & 16.0 & \textbf{20.0} & 24.0 & 28.0 & 32.0 & 36.0 & 40.0 \\
\midrule
最优 & R & R & R & R & R/A & A & A & A & A & A \\
\bottomrule
\end{tabular}
\end{center}

\begin{warnbox}[title={v1.0 的“三段结构”要改口，但不能说 SAFE 永远出局}]
v1.0 里 SAFE 恒为 25（$3$ 食物/人 $\times4=12$，加 10 英里），不随 $s$ 变，
所以 $s\ge8$ 时 SAFE 反超。v1.1 的 SAFE \textbf{已经随 $s$ 缩放}，那个结论作废。

但如果只按 v1.0 的形式代入当前价格估算：
\[
C(\text{SAFE})\approx 3m\lambda_f+10=12+10=22
\]
与 $C(\text{AVOID})=4s$ 比较，$4s>22\Rightarrow s\ge6$ 时 \textbf{SAFE 反而更优}。
也就是说，\textbf{约 $s\ge6$ 这一段要留给 SAFE}。

\textbf{但这是下限。}官方说 SAFE 现在“随严重度缩放”，实际成本一定高于 22，
所以反超点会\textbf{右移甚至消失}——具体取决于未公布的新公式。

\noindent 稳妥做法：$s\ge6$ 时把 SAFE 和 AVOID 都算一遍再选，
不要写死成“一律 AVOID”。这一条放进待确认清单。
\end{warnbox}

\subsubsection*{ATTACK}

\begin{align*}
C(\text{SAFE})&=(20+s)\lambda_m+10\lambda_f=20+0.5s\\
C(\text{RISKY})&=m s\lambda_h d=4s\\
C(\text{AVOID})&=15
\end{align*}

$4s<15\Rightarrow s\le3$：\textbf{RISKY}；$s\ge4$：\textbf{AVOID}。SAFE 永远最差。

\subsubsection*{THEFT}

\begin{align*}
C(\text{SAFE})&=(5+s)\lambda_f=5+s\\
C(\text{RISKY})&=0.5(5m\lambda_h d)+0.5(20\lambda_f+20\lambda_m)=25\\
C(\text{AVOID})&=15\lambda_m=7.5
\end{align*}

$5+s<7.5\Rightarrow s\le2$：\textbf{SAFE}；$s\ge3$：\textbf{AVOID}。
注意这与 v1.0 相反（v1.0 里 SAFE 恒为 10 食物，永远输）。

\subsubsection*{BAD\_TRAIL}

\begin{align*}
C(\text{SAFE})&=6+s\quad(\text{若仍扣 5 食物则}+5)\\
C(\text{RISKY})&=\tfrac{s}{10}\big(m s\lambda_h d+0.15\lambda_o\big)-\big(1-\tfrac{s}{10}\big)\cdot5
=0.4s^2+1.4s-5\\
C(\text{AVOID})&=20
\end{align*}

$C(\text{RISKY})=C(\text{SAFE})$ 在 $s\approx4.5$ 处相交：
\textbf{$s\le4$ RISKY，$s\ge5$ SAFE}。
$s=1$ 时 RISKY 是\textbf{负代价}（$-3.2$）——大概率白赚 5 英里。

\subsubsection*{HUNT\_OPP}

\begin{align*}
C(\text{SAFE})&=-20\lambda_f+2\lambda_a=-17\\
C(\text{RISKY})&=-50\lambda_f+4\lambda_a-\tfrac{s}{10}(5m\lambda_h d)=-44-2.0s\\
C(\text{AVOID})&=0
\end{align*}

两者都净赚，RISKY 赚得更多。\textbf{只要弹药够（$\ge4$）且车厢没满，RISKY 恒胜。}
这是全场收益最高的事件，值得为它预留弹药。

\subsubsection*{TRADER}

\[
\text{买入条件}:\ P_{\text{food}}<\frac{\lambda_f}{\lambda_m}=2.0,\qquad
\text{卖出条件}:\ P_{\text{sell}}>\lambda_x
\]

v1.1 的 $\lambda_f/\lambda_m=2.0$：\$1.4 买、\$2.0 打平、\$2.8 不买。
砍价期望 $=q\cdot0.3\cdot P\cdot n$，$q=0.5+(s-5)\times0.04$。

\subsection{汇总表（$m=4$，含 Doctor）}

\begin{center}
\small
\begin{longtable}{@{}p{2.8cm}p{3.2cm}p{3.4cm}p{2.2cm}p{3.2cm}@{}}
\toprule
\textbf{事件} & \textbf{SAFE} & \textbf{RISKY} & \textbf{AVOID} & \textbf{结论} \\
\midrule
\endhead
RIVER & $10+0.5s$ & $6s+12$ & 5 & AVOID 恒胜 \\
STORM & $4s$ & $4s+9$ & 10 & $s\le2$ S，$s\ge3$ A；RISKY 被支配 \\
BREAKDOWN & 15（Mechanic: 0） & $3.5s$ & 15 & 有 Mechanic 就 SAFE \\
ILLNESS & 未公布（$\ge22$） & $0.4s(s+5)$ & $4s$ & $s\le4$ R；$s=5$ R/A 打平；$s\ge6$ 须比较 SAFE 与 AVOID \\
ATTACK & $20+0.5s$ & $4s$ & 15 & $s\le3$ R，$s\ge4$ A \\
THEFT & $5+s$ & 25 & 7.5 & $s\le2$ S，$s\ge3$ A \\
BAD\_TRAIL & $6+s$ & $0.4s^2+1.4s-5$ & 20 & $s\le4$ R，$s\ge5$ S \\
HUNT\_OPP & $-17$ & $-44-2.0s$ & 0 & 有弹药就 RISKY \\
\bottomrule
\end{longtable}
\end{center}

\subsection{三条覆盖规则}

公式之外，还有三个非线性因素必须强制处理：

\begin{center}
\small
\begin{longtable}{@{}p{0.7cm}p{5.0cm}p{4.0cm}p{4.2cm}@{}}
\toprule
\textbf{\#} & \textbf{条件} & \textbf{强制选择} & \textbf{原因} \\
\midrule
\endhead
R1 & 任一成员 HP $\in[s,\ s+5)$ & AVOID & RISKY 伤害 $-(s+5)$ 会直接打死 \\
R2 & \texttt{oxen} $=4$ 且 $(2000-D)/15>T_{\text{left}}$ & 禁止 RISKY & 掉到 3 头后，剩余回合跑不完剩余里程 \\
R3 & food $<3m$ & 禁止 SAFE（食物类事件） & 会触发饥饿，$5\to23$ HP/人 \\
\bottomrule
\end{longtable}
\end{center}

\begin{keybox}[title={R2 是 v1.1 新增的，也是最重要的一条}]
\textbf{前提要算清楚：}3 头牛 = 15 英里/回合。如果此时剩余里程 $2000-D$ 除以 15
超过剩余回合 $T_{\text{left}}$，那掉牛就是\textbf{必败}，不是“慢一点”。

所以判据是条件式，不是“4 头牛一律不冒险”：

\begin{center}
\begin{tabular}{@{}lcc@{}}
\toprule
\textbf{时点} & \textbf{掉到 3 牛后的所需回合} & \textbf{是否致命} \\
\midrule
第 20 回合（剩 1800 英里 / 100 回合） & $1800/15=120$ & \textbf{是}，必须严禁 \\
第 60 回合（剩 800 英里 / 60 回合） & $800/15=54$ & 否，可以承担 \\
第 90 回合（剩 200 英里 / 30 回合） & $200/15=14$ & 否，完全无所谓 \\
\bottomrule
\end{tabular}
\end{center}

\textbf{补充：牛是可以在交易站补买的}（TRADER SAFE，\$80/头），所以“掉牛=永久出局”是过度概括。
真正不可逆的只有两件事：\textbf{死亡}和\textbf{回合耗尽}。

结论：\textbf{前半程把 R2 当铁律，后半程按上表算。}
\end{keybox}

%=====================================================================
\section{决策面 D：动作分配（v1.1）}

\subsection{动作的英里当量}

\begin{center}
\begin{tabular}{@{}lccc@{}}
\toprule
\textbf{动作} & \textbf{直接产出} & \textbf{英里当量} & \textbf{约束} \\
\midrule
HUNT（Hunter） & $+10$ 食物，$-2$ 弹药 & $10-3=7.0$ & 弹药 $\ge2$，food $<500$ \\
HUNT（非 Hunter） & $+7$ 食物，$-2$ 弹药 & $7-3=4.0$ & 同上 \\
HEAL（Doctor） & $+5$ HP 给最伤员，$-3$ 食物 & $5\times2.0-3=7.0$ & 有人受伤，且不挨饿 \\
DEFEND & 全队 $+3$ HP $=+12$ HP & $12\times2.0\times0.5=12.0$ & 有事件，全局只算一次 \\
SCOUT（Scout） & $+6$ 英里 & $6.0$ & 无 \\
MOVE & $+5$ 英里，$-1$ HP & $5-2.0=3.0$ & 无 \\
HEAL（非 Doctor） & $+3$ HP，$-3$ 食物 & $3\times2.0-3=3.0$ & 同上 \\
REST & $+2$ HP & $4.0$ & 同上 \\
\bottomrule
\end{tabular}
\end{center}

\begin{keybox}[title={v1.1 的动作排序变了吗}]
\textbf{DEFEND 从第四名升到第一。}因为事件变多、伤害变贵，一个人防一次值 12 英里。

\textbf{第二名是并列的 HUNT（Hunter）和 Doctor 的 HEAL，各 7.0。}
这意味着：\textbf{在 v1.1 里，Doctor 治疗和 Hunter 打猎一样值钱}——
因为它同时还是减伤光环。

\textbf{MOVE 掉到只剩 3.0}，比 SCOUT（6.0）差一半。
\textbf{非 Doctor 的 HEAL 也是 3.0}，除非濒死否则不值得做。
\end{keybox}

\subsection{分配算法}

\begin{enumerate}[leftmargin=1.8em]
  \item 若有事件（\texttt{isThreatened()}）：指定\textbf{一个}成员 DEFEND（通常 Mechanic）。
  \item 若 runway $<$ 阈值：Hunter 优先 HUNT。
  \item 若有人 HP $<$ 阈值：Doctor 优先 HEAL。
  \item Scout 默认 SCOUT。
  \item 剩余动作：比较 HUNT 与 MOVE 的边际价值，取大者。
\end{enumerate}

\begin{warnbox}[title={两个必须写死的约束}]
\textbf{1. DEFEND 全局只算一次。}第二个防守者的边际收益是 0。
\textbf{2. 挨饿回合 HEAL/REST 无效。}此时它们的英里当量是 0，必须改成 HUNT。
\end{warnbox}

%=====================================================================
\section{四层如何协调：\texttt{Plans} 的契约}

\begin{center}
\small
\begin{tabular}{@{}p{3.2cm}p{4.6cm}p{6.8cm}@{}}
\toprule
\textbf{阶段} & \textbf{谁写} & \textbf{写什么} \\
\midrule
1. \texttt{decide()} & \texttt{TeamStrategy} & 完整状态快照 + 两个开关 \\
4. \texttt{resolve()} & 只读 & 用快照弥补“看不见车” \\
5. \texttt{act()} & 只读 & 用快照统一口径，但成员另有 \texttt{Situation} \\
\bottomrule
\end{tabular}
\end{center}

\begin{center}
\begin{tabular}{@{}p{4.0cm}p{2.4cm}p{8.2cm}@{}}
\toprule
\textbf{字段} & \textbf{类型} & \textbf{用途} \\
\midrule
\texttt{foodRunway} & int & 食物还能撑几回合；派生 \texttt{foodScarce} \\
\texttt{livingMembers} & int & 所有公式里的 $m$ \\
\texttt{minHealth} & int & R1 判断“最伤员” \\
\texttt{averageHealth} & int & 决定口粮档位 \\
\texttt{oxen} & int & \textbf{R2 的关键输入} \\
\texttt{food}、\texttt{ammo} 等资源 & int & 交易判断 \\
\texttt{turnsLeft} & int & 判断 MOVE 的价值是否上升 \\
\texttt{foodScarce} / \texttt{turnsTight} & boolean & 供 Resolver 用 \\
\bottomrule
\end{tabular}
\end{center}

\begin{keybox}[title={v1.1 新增：\texttt{oxen} 必须进 \texttt{Plans}}]
R2 规定“牛等于 4 时禁止有死牛风险的 RISKY”，
而 \texttt{Resolver} 看不到车队——所以 \textbf{\texttt{oxen} 必须由策略层写进 \texttt{Plans}}。
这和 \texttt{minHealth} 一样，是共享上下文存在的硬理由。
\end{keybox}

%=====================================================================
\section{完整判定顺序（伪代码）}

\begin{verbatim}
// ---- decide(state) ----
plans.turn          = state.getTurn()
plans.livingMembers = state.aliveMembers()
plans.oxen          = amountOf(OXEN)
plans.minHealth     = min health over living members
plans.averageHealth = mean health over living members
plans.turnsLeft     = 120 - turn

runway           = food / max(1, 3*living + 2*oxen)
plans.foodRunway = runway
plans.foodScarce = runway < 3
plans.turnsTight = (2000 - distance) / turnsLeft > 20

ration = runway < 3                      ? BARE
       : (avgHealth < 50 && runway > 5)  ? FULL
       : MEAGER
pace   = turnsTight ? GRUELING : NORMAL
return new Directive(pace, ration)

// ---- resolve(event) ----
s = event.severity();  m = plans.livingMembers
if (event.isMarket()) return tradeRule(event)

base = table(event.kind(), s, m)          // 7.2 的公式

// R2: risky ox loss is fatal only if 3 oxen can't finish in time
boolean oxLossFatal = plans.oxen == 4
    && (2000 - distance) / 15 > plans.turnsLeft;
if (oxLossFatal && base == RISKY
    && kind in {RIVER, STORM, BAD_TRAIL}) base = SAFE

if (kind == ILLNESS && minHealth in [s, s+5)) base = AVOID    // R1
if (kind == ILLNESS && minHealth < s)         base = RISKY    // R1b
if (kind in {ILLNESS, BAD_TRAIL} && food < 3*m) base = AVOID  // R3

return Resolution.of(base)

// ---- act(situation) ----
if (starving)                       return HUNT
if (isThreatened() && amDefender)   return DEFEND
if (amHunter && runway < LOW_FOOD)  return HUNT
if (amDoctor && minHealth < LOW_HP) return HEAL
if (amScout)                        return SCOUT
return MOVE
\end{verbatim}

%=====================================================================
\section{校准：用回放修正影子价格}

模型的价值取决于 $\lambda$ 准不准。三个校准实验，按顺序做：

\begin{center}
\small
\begin{longtable}{@{}p{0.8cm}p{4.4cm}p{4.6cm}p{4.6cm}@{}}
\toprule
\textbf{\#} & \textbf{实验} & \textbf{怎么做} & \textbf{看什么} \\
\midrule
\endhead
E1 & 食物价格 $\lambda_f$ & 固定其他变量，只改口粮 BARE vs MEAGER & 距离差 / 食物差 \\
E2 & HP 价格 $\lambda_h$ & 固定采购，只改治疗频率 & 血线差与距离差的比值 \\
E3 & 牛的价值 $\lambda_o$ & 4 牛 vs 5 牛各跑一次 & 到达率（不是分数） \\
\bottomrule
\end{longtable}
\end{center}

\begin{goodbox}[title={校准的收敛标准}]
三个种子都到 2000 英里、无死亡、饥饿回合 $\le2$ 时，\textbf{停止调参}。
继续调只会开始过拟合这三个种子。
\end{goodbox}

%=====================================================================
\section{模型的已知局限}

\begin{enumerate}[leftmargin=1.8em]
  \item \textbf{线性化。}HP 的真实价值是凸的，线性 $\lambda_h$ 在濒死时会低估死亡风险。
  \item \textbf{单步近视。}每回合重算，看不见“今天的选择对 10 回合后的影响”。
  \item \textbf{两个关键公式缺失。}事件概率、严重度、ILLNESS SAFE 的 v1.1 公式都没公布。
        本文用 v1.0 的结构 + v1.1 的已知值推算，\textbf{结论可能随官方公式公布而偏移}。
  \item \textbf{三种子过拟合。}所有参数都是对 3 个固定种子的拟合。
        \textbf{规则要基于机制，不要基于具体回合数。}
  \item \textbf{$\lambda_o$ 不是线性价格。}牛的价值是断崖式的（4 头 vs 3 头），
        用 60 这个数字只是把它压成一个可比较的量；真正的判据是 R2 的回合预算。
  \item \textbf{$\lambda_h$ 取 2.0 是估值，不是事实。}如果实测发现队伍血线普遍偏高，
        说明真实 $\lambda_h$ 更低，可以更激进地选 RISKY。它的合理区间约在 $[1.5,3.0]$，
        会让部分阈值移动 $\pm1$ 个 $s$，但\textbf{不改变“哪个姿态被支配”这个结论}。
\end{enumerate}

%=====================================================================
\section{待确认}

\begin{center}
\begin{longtable}{@{}p{0.8cm}p{7.4cm}p{6.4cm}@{}}
\toprule
\textbf{\#} & \textbf{事项} & \textbf{影响} \\
\midrule
\endhead
1 & ILLNESS 的 SAFE 新公式 & 决定 SAFE 是否还有可能胜出 \\
2 & 事件概率公式（已知 20\% $\to$ 56\%） & 事件预算与期望计算 \\
3 & 严重度新公式 & 分段阈值的位置 \\
4 & HUNT 产出是否仍为 7/10 & 食物预算的整体比例 \\
5 & BAD\_TRAIL 的 SAFE 是否仍扣 5 食物 & 该事件的 $s\ge5$ 区间 \\
6 & Doctor 光环“减半”的取整方式 & 边界伤害计算 \\
7 & cargo 是否封顶 1000 & 后期是否值得囤货 \\
\bottomrule
\end{longtable}
\end{center}

\vfill
\begin{center}
\small\color{muted}
配套文档：\texttt{01} 规则 \quad \texttt{02} 评分 \quad \texttt{03} 项目方案 \quad
\texttt{05} 分工提案
\end{center}

\end{document}
